\documentclass[border=12pt]{standalone}
\usepackage[utf8]{inputenc}
\usepackage[T1]{fontenc}
\usepackage{amsmath}
\usepackage[siunitx, RPvoltages, americanresistors]{circuitikz}
\sisetup{output-decimal-marker = {,}}

\begin{document}
\begin{circuitikz}[scale=0.92, transform shape, font=\small]

    % ==================== ETAPA 1 (Q1 = 0.541, G1 = 2) ====================
    \draw (0,0) node[left] {$V_{\text{in}}$} to[short, o-] (0.6,0)
          to[R, l=$R_1$, a={\footnotesize $3{,}9\,\text{k}\Omega$}] (2.8,0) coordinate (n1)
          to[R, l=$R_2$, a={\footnotesize $24\,\text{k}\Omega$}] (5.0,0) coordinate (n2)
          -- (6.4,0);

    % C1 a masa (bien separado a la izquierda del AO)
    \draw (n2) to[C, l_={\footnotesize $C_1 = 2{,}2\,\text{nF}$}, *-] (5.0,-2.1) node[ground]{};

    % Op-Amp U1A (TL084: + Pin 3, - Pin 2, Out Pin 1)
    \node[op amp, noinv input up, anchor=+] (OA1) at (6.4,0) {};
    \node[font=\footnotesize\bfseries] at (OA1.center) {U1A};
    \node[above left, font=\scriptsize\bfseries, color=blue!75!black] at (OA1.+) {3};
    \node[below left, font=\scriptsize\bfseries, color=blue!75!black] at (OA1.-) {2};
    \node[above right, font=\scriptsize\bfseries, color=blue!75!black] at (OA1.out) {1};

    % Alimentacion compacta en U1A
    \draw (OA1.up) -- ++(0,0.28) node[above, font=\scriptsize]{$+V_{CC}$ (4)};
    \draw (OA1.down) -- ++(0,-0.22) node[below, font=\scriptsize]{$-V_{EE}$ (11)};

    % Salida Etapa 1 y Realimentacion positiva C2 por arriba (y = +2.5)
    \draw (OA1.out) -- ++(0.75,0) coordinate (out1);
    \draw (out1) to[short, *-] (out1 |- 0,2.5)
          to[C, l_=$C_2$, a^={\footnotesize $3{,}3\,\text{nF}$}] (n1 |- 0,2.5)
          to[short, -*] (n1);

    % Red de realimentacion negativa G1 = 2 (Ra1, Rb1) bien abajo (y = -2.4)
    \draw (OA1.-) -- (5.85, 0 |- OA1.-) -- (5.85,-2.45) coordinate (inv1b);
    \draw (inv1b) to[R, l_={\footnotesize $R_{a1}=10\,\text{k}\Omega$}, *-] (5.85,-4.35) node[ground]{};
    \draw (inv1b) to[R, l=$R_{b1}$, a={\footnotesize $10\,\text{k}\Omega$}] (out1 |- inv1b)
          to[short, -*] (out1);

    \node[above, font=\footnotesize\bfseries, color=purple!80!black] at ([xshift=0.82cm]out1) {$V_{\text{etapa1}}$};

    % ==================== ETAPA 2 (Q2 = 1.307, G2 = 2) ====================
    \draw (out1) -- ++(1.65,0) coordinate (in2)
          to[R, l=$R_3$, a={\footnotesize $2{,}0\,\text{k}\Omega$}] ++(2.2,0) coordinate (n3)
          to[R, l=$R_4$, a={\footnotesize $10\,\text{k}\Omega$}] ++(2.2,0) coordinate (n4)
          -- ++(1.4,0) coordinate (plus2);

    % C3 a masa
    \draw (n4) to[C, l_={\footnotesize $C_3 = 3{,}3\,\text{nF}$}, *-] ([yshift=-2.1cm]n4) node[ground]{};

    % Op-Amp U1B (TL084: + Pin 5, - Pin 6, Out Pin 7)
    \node[op amp, noinv input up, anchor=+] (OA2) at (plus2) {};
    \node[font=\footnotesize\bfseries] at (OA2.center) {U1B};
    \node[above left, font=\scriptsize\bfseries, color=blue!75!black] at (OA2.+) {5};
    \node[below left, font=\scriptsize\bfseries, color=blue!75!black] at (OA2.-) {6};
    \node[above right, font=\scriptsize\bfseries, color=blue!75!black] at (OA2.out) {7};

    % Salida Etapa 2 y Realimentacion positiva C4 por arriba
    \draw (OA2.out) -- ++(0.75,0) coordinate (out2);
    \draw (out2) to[short, *-] ([yshift=2.5cm]out2) coordinate (top2)
          to[C, l_=$C_4$, a^={\footnotesize $10\,\text{nF}$}] (n3 |- top2)
          to[short, -*] (n3);

    % Red de realimentacion negativa G2 = 2 (Ra2, Rb2)
    \coordinate (inv2x) at ([xshift=-0.55cm]OA2.-);
    \draw (OA2.-) -- (inv2x) -- ([yshift=-1.96cm]inv2x) coordinate (inv2b);
    \draw (inv2b) to[R, l_={\footnotesize $R_{a2}=10\,\text{k}\Omega$}, *-] ([yshift=-1.9cm]inv2b) node[ground]{};
    \draw (inv2b) to[R, l=$R_{b2}$, a={\footnotesize $10\,\text{k}\Omega$}] (out2 |- inv2b)
          to[short, -*] (out2);

    % Salida final
    \draw (out2) to[short, -o] ++(1.1,0) node[right] {$V_{\text{out}}$ \footnotesize ($G_0 = 4\,\text{V/V} = 12{,}04\,\text{dB}$)};

    % ==================== RECUADROS DE ETAPAS ====================
    \draw[dashed, rounded corners=4pt, gray!75!black, line width=0.6pt] (0.75,-5.25) rectangle (9.80, 3.65);
    \node[below, font=\footnotesize\bfseries, color=black!80] at (5.27,-5.32)
        {Celda 1 (Sallen-Key): $f_{01}=6{,}11\,\text{kHz},\; Q_1=0{,}537,\; G_1=2\,\text{V/V}$};

    \draw[dashed, rounded corners=4pt, gray!75!black, line width=0.6pt] (10.75,-6.15) rectangle (20.15, 3.15);
    \node[below, font=\footnotesize\bfseries, color=black!80] at (15.45,-6.22)
        {Celda 2 (Sallen-Key): $f_{02}=6{,}20\,\text{kHz},\; Q_2=1{,}311,\; G_2=2\,\text{V/V}$};

    % ==================== NOTA DE DESACOPLE TL084 ====================
    \node[draw=gray!70!black, fill=yellow!8, rounded corners=3pt, align=left, font=\scriptsize, anchor=north west] at (0.75, -6.90) {
        \textbf{Integrado TL084 (DIP-14) — Notas para Layout Placa 2:} Alimentación simétrica $+V_{CC}=+15\,\text{V}$ (o $+12\,\text{V}$) en \textbf{Pin 4} y $-V_{EE}=-15\,\text{V}$ (o $-12\,\text{V}$) en \textbf{Pin 11}.\\
        Colocar $2 \times C_{\text{desacople}} = 100\,\text{nF}$ cerámicos lo más cerca posible del integrado (uno entre Pin 4 y GND, otro entre Pin 11 y GND).\\
        Las secciones \textbf{U1C} (pines 10+, 9$-$, 8 out) y \textbf{U1D} (pines 12+, 13$-$, 14 out) quedan libres para el Pasa Altos de orden 4 (o como seguidor con entrada $+$ a GND si usan IC separado).
    };

\end{circuitikz}
\end{document}
